% ═══════════════════════════════════════════════════════════════
%  ch01.tex — فصل ۱: مفهوم‌شناسی مداخله‌ی خارجی و نجات‌بخشی
%  (ادامه — از محل قطع‌شده‌ی پیوستار مداخله)
% ═══════════════════════════════════════════════════════════════

\chapter{مفهوم‌شناسی مداخله‌ی خارجی و نجات‌بخشی}
\label{ch:concepts}

\begin{tcolorbox}[quotebox, title={}]
«هیچ ملتی نمی‌تواند آزاد باشد و در عین حال ملت‌های دیگر را در بند نگاه دارد.»\\
\hfill --- فریدریش انگلس\src{\lr{Engels, 1847}}
\end{tcolorbox}

% ────────────────────────────────────────────────────────────────
\section{درآمد}
\label{sec:ch01-intro}

مداخله‌ی خارجی در امور داخلی ملت‌ها یکی از مناقشه‌برانگیزترین پدیده‌های
نظام بین‌المللی است. از سویی، اصل \concept{حاکمیت ملی}%
\footnote{\lat{National Sovereignty}}
--- که از معاهده‌ی وستفالی (۱۶۴۸ م.) به‌عنوان ستون فقرات نظم بین‌المللی
شناخته می‌شود --- هرگونه دخالت بیرونی در سرنوشت یک واحد سیاسی را ناموجه
می‌شمارد. از سوی دیگر، تجربه‌ی تاریخی آکنده از نمونه‌هایی است که در آن‌ها
گروه‌هایی در درون یک واحد سیاسی، خود دست به طلب کمک از نیروهای خارجی
زده‌اند --- خواه با انگیزه‌ی رهایی از ستم، خواه با محاسبه‌های قدرت، و
خواه از سر ناامیدی.

\thinker{مایکل والزر}%
\footnote{\lat{Michael Walzer (1935–)},
فیلسوف سیاسی آمریکایی و نویسنده‌ی \lat{\textit{Just and Unjust Wars}} (۱۹۷۷).}
در اثر تأثیرگذار خود استدلال می‌کند که مداخله‌ی نظامی تنها در موارد بسیار
محدود --- از جمله نسل‌کشی و بردگی سیستماتیک --- قابل توجیه اخلاقی
است\src{\lr{Walzer, 1977, pp.\,86--108}}.
اما واقعیت تاریخی بسی پیچیده‌تر از این چارچوب هنجاری است. در بسیاری از موارد، مرز میان «دعوت برای آزادسازی» و «توجیه تجاوز» چنان نازک بوده که تنها با فاصله‌ی زمانی و بررسی نتایج بلندمدت می‌توان قضاوتی نسبتاً موجه عرضه کرد.

هدف این فصل، تبیین دقیق مفاهیم کلیدی پژوهش حاضر است: \concept{مداخله‌ی
خارجی}، \concept{نجات‌بخشی}، \concept{نجات‌طلبی}، \concept{دعوت از
بیگانه}، و طیف گسترده‌ای از مفاهیم مرتبط که بدون روشن‌سازی آن‌ها، تحلیل
موردی فصل‌های آتی امکان‌پذیر نخواهد بود.


% ────────────────────────────────────────────────────────────────
\section{مداخله‌ی خارجی: تعریف و طیف مفهومی}
\label{sec:ch01-intervention-def}

\begin{tcolorbox}[defbox, title={تعریف: مداخله‌ی خارجی}]
\concept{مداخله‌ی خارجی}%
\footnote{\lat{Foreign Intervention}}
به هرگونه کنش عمدی یک بازیگر خارجی (دولت، سازمان بین‌المللی، نیروی
نظامی غیردولتی، یا حتی بازیگران اقتصادی و فرهنگی) اطلاق می‌شود که
با هدف تغییر رفتار، ساختار سیاسی، یا توازن قدرت در درون یک واحد
سیاسی دیگر صورت می‌گیرد، فارغ از اینکه این کنش با دعوت داخلی
همراه باشد یا بدون آن.
\end{tcolorbox}

این تعریف عامدانه گسترده طراحی شده است تا طیف کاملی از پدیده‌ها را پوشش
دهد. \thinker{جیمز روزنا}%
\footnote{\lat{James N.\,Rosenau (1924--2011)},
نظریه‌پرداز روابط بین‌الملل.}
در طبقه‌بندی سه‌گانه‌ی خود میان مداخله‌ی \concept{دیپلماتیک}،
\concept{اقتصادی} و \concept{نظامی} تمایز قائل می‌شود%
\src{\lr{Rosenau, 1969, pp.\,152--167}}.\\
اما طبقه‌بندی‌های متأخرتر --- به‌ویژه پس از جنگ سرد --- ابعاد \concept{فرهنگی-ایدئولوژیک}، \concept{سایبری} و \concept{اطلاعاتی-رسانه‌ای} را نیز افزوده‌اند.

\subsection{طیف مداخله: از نرم تا سخت}
\label{subsec:ch01-spectrum}

برای درک بهتر، طیف مداخله را می‌توان در یک پیوستار ترسیم کرد:

\begin{figure}[htbp]
\centering
\begin{tikzpicture}[
  >=Stealth,
  boxstyle/.style={
    draw, rounded corners=3pt, minimum height=1.2cm,
    text width=2.6cm, align=center, font=\small\bfseries
  }
]
% محور اصلی
\draw[->, very thick, PrimaryDeep]
  (0,0) -- (16,0)
  node[below, font=\footnotesize, xshift=-1cm] {\rl{شدت مداخله}};

% گره‌ها
\node[boxstyle, fill=BgTeal, text=AccentTeal] (a) at (1.5,1.5)
  {\rl{فرهنگی}\\[-2pt]{\footnotesize\rl{میسیونری}}};
\node[boxstyle, fill=BgWarm, text=AccentAmber] (b) at (4.5,1.5)
  {\rl{دیپلماتیک}\\[-2pt]{\footnotesize\rl{فشار سیاسی}}};
\node[boxstyle, fill=BgCool, text=PrimaryMid] (c) at (7.5,1.5)
  {\rl{اقتصادی}\\[-2pt]{\footnotesize\rl{تحریم/کمک}}};
\node[boxstyle, fill=yellow!15, text=EraContemp] (d) at (10.5,1.5)
  {\rl{اطلاعاتی}\\[-2pt]{\footnotesize\rl{سایبری/رسانه}}};
\node[boxstyle, fill=red!10, text=AccentRed] (e) at (13.5,1.5)
  {\rl{نظامی}\\[-2pt]{\footnotesize\rl{حمله/اشغال}}};

% خطوط اتصال به محور
\foreach \nd in {a,b,c,d,e}{%
  \draw[thick, NeutralMid] (\nd.south) -- (\nd.south |- 0,0);%
}%

% برچسب‌های زیر محور
\node[below, font=\scriptsize, text=AccentTeal] at (1.5,-0.3) {\rl{نرم}};
\node[below, font=\scriptsize, text=AccentRed] at (13.5,-0.3) {\rl{سخت}};

% تیتر
\node[above, font=\normalsize\bfseries, text=PrimaryDeep]
  at (8,3) {\rl{پیوستار شدت مداخله‌ی خارجی}};
\end{tikzpicture}
\caption{طیف مداخله‌ی خارجی از نرم (فرهنگی) تا سخت (نظامی)}
\label{fig:ch01-spectrum}
\end{figure}

نکته‌ی مهم آن است که این طیف یک‌بعدی نیست. یک مداخله‌ی ظاهراً «نرم» ---
مانند فعالیت‌های میسیونری --- ممکن است تأثیرات ساختاری عمیق‌تری از یک
تهاجم نظامی کوتاه‌مدت داشته باشد. \thinker{جوزف نای}%
\footnote{\lat{Joseph S.\,Nye Jr.~(1937–)},
نظریه‌پرداز «قدرت نرم» \lat{(Soft Power)}.}
نشان داده است که قدرت نرم --- از طریق فرهنگ، ارزش‌ها و نهادها ---
می‌تواند «ترجیحات دیگران را شکل دهد» بدون آنکه ظاهراً اجباری در
کار باشد\src{\lr{Nye, 2004, pp.\,5--11}}.


\subsection{مداخله‌ی دعوت‌شده در برابر مداخله‌ی تحمیلی}
\label{subsec:ch01-invited-vs-imposed}

تمایز بنیادینی که در سرتاسر این پژوهش محوری خواهد بود، تفکیک میان
\concept{مداخله‌ی دعوت‌شده}\footnote{\lat{Intervention by Invitation}}
و \concept{مداخله‌ی تحمیلی}\footnote{\lat{Imposed Intervention}} است.

\begin{tcolorbox}[comparebox, title={\rl{مقایسه: مداخله‌ی دعوت‌شده و تحمیلی}}]
\begin{longtable}{>{\bfseries\color{PrimaryMid}}p{3cm}
  p{5.5cm} p{5.5cm}}
\toprule
\textbf{معیار} & \textbf{مداخله‌ی دعوت‌شده} & \textbf{مداخله‌ی تحمیلی} \\
\midrule
\endfirsthead
\toprule
\textbf{معیار} & \textbf{مداخله‌ی دعوت‌شده} & \textbf{مداخله‌ی تحمیلی} \\
\midrule
\endhead
منشأ تصمیم &
بازیگر(ان) داخلی درخواست می‌کنند &
بازیگر خارجی یک‌جانبه اقدام می‌کند \\
\midrule
مشروعیت ادعایی &
رضایت داخلی (هرچند جزئی) &
ضرورت بین‌المللی یا توجیه اخلاقی \\
\midrule
نمونه‌ی تاریخی &
ویلیام اورانج در انگلستان (۱۶۸۸) &
حمله‌ی آمریکا به عراق (۲۰۰۳) \\
\midrule
مسئله‌ی کلیدی &
نمایندگی دعوت‌کننده: آیا نماینده‌ی ملت است؟ &
حق مداخله: آیا بازیگر خارجی چنین حقی دارد؟ \\
\bottomrule
\end{longtable}
\end{tcolorbox}

\thinker{لوئیز داسپرمون}%
\footnote{\lat{Jean d'Aspremont},
حقوقدان بین‌المللی بلژیکی.}
در پژوهش مهم خود درباره‌ی مداخله به دعوت، نشان می‌دهد که حتی
مفهوم «دعوت» خود مناقشه‌آمیز است: چه کسی حق دعوت دارد؟ دولت
مستقر یا مخالفان آن؟ اکثریت یا اقلیت ستمدیده؟%
\src{\lr{d'Aspremont, 2015, pp.\,24--48}}


% ────────────────────────────────────────────────────────────────
\section{نجات‌بخشی و نجات‌طلبی: دو سوی یک سکه}
\label{sec:ch01-rescue}

\begin{tcolorbox}[defbox, title={تعریف: نجات‌بخشی}]
\concept{نجات‌بخشی}%
\footnote{\lat{Deliverance / Rescue / Liberation}}
وضعیتی است که در آن یک بازیگر خارجی --- با ادعا یا واقعیت --- نقش
رهاننده‌ی گروهی از مردم از وضعیت ستم، بحران یا فروپاشی را ایفا
می‌کند. نجات‌بخشی ممکن است خودخوانده باشد (ادعای بازیگر خارجی)
یا خواسته‌شده (پاسخ به درخواست داخلی).
\end{tcolorbox}

\begin{tcolorbox}[defbox, title={تعریف: نجات‌طلبی}]
\concept{نجات‌طلبی}%
\footnote{\lat{Salvation-Seeking / Appeal for External Deliverance}}
وضعیتی است که در آن یک گروه داخلی --- اعم از نخبگان سیاسی، اقلیت‌های
قومی-مذهبی، مرزنشینان، یا حتی بخش‌هایی از عوام --- به‌طور فعال
در جستجوی یک ناجی خارجی برمی‌آیند و از نیروی بیرونی برای رهایی
از وضعیت موجود مدد می‌طلبند.
\end{tcolorbox}

تفاوت ظریف ولی تعیین‌کننده‌ای میان این دو مفهوم وجود دارد. در
نجات‌بخشی، کانون عاملیت در بیرون است؛ در نجات‌طلبی، آغازگر درون
است اما ابزار در بیرون. با این حال، در واقعیت تاریخی این دو اغلب
چنان درهم‌تنیده‌اند که تفکیکشان نیازمند تحلیل لایه‌ای است ---
تحلیلی که در فصل بعد به تفصیل خواهیم پرداخت.


% ────────────────────────────────────────────────────────────────
\section{عاملیت، بازنمایی و مسئله‌ی نمایندگی}
\label{sec:ch01-agency}

یکی از پرسش‌های بنیادین در هر مداخله‌ی «دعوت‌شده» آن است که
\concept{دعوت‌کننده تا چه حد نماینده‌ی مردم است}. این پرسش را
\thinker{گایاتری اسپیواک}\footnote{\lat{Gayatri Chakravorty Spivak (1942–)},
نظریه‌پرداز پسااستعماری هندی‌تبار.}
در مقاله‌ی مشهور خود به شکل تیزتری مطرح کرده است: «آیا فرودست
می‌تواند سخن بگوید؟»\src{\lr{Spivak, 1988, pp.\,271--313}}

مسئله آن است که در بسیاری از موارد تاریخی:
\begin{enumerate}[label=\textbf{\arabic*.}, nosep, rightmargin=1em]
  \item دعوت‌کننده یک نخبه‌ی سیاسی بوده، نه مردم عادی (مانند شاهزادگان
    ساسانی که از اعراب مدد خواستند)؛
  \item دعوت‌کننده یک اقلیت ناراضی بوده که نه لزوماً نماینده‌ی اکثریت
    بوده و نه منافعش با منافع ملی همسو (مانند شیخ خزعل و بریتانیا)؛
  \item دعوت‌کننده اساساً ساختگی بوده و قدرت خارجی برای مشروعیت‌بخشی
    به مداخله‌ی خود، یک «متقاضی» دست‌نشانده تراشیده است.
\end{enumerate}

\subsection{مسئله‌ی اطلاعات ناقص و تصمیم‌گیری در بحران}
\label{subsec:ch01-info}

\thinker{هربرت سایمون}%
\footnote{\lat{Herbert A.\,Simon (1916--2001)},
برنده‌ی جایزه‌ی نوبل اقتصاد (۱۹۷۸) و نظریه‌پرداز «عقلانیت محدود».}
نشان داده است که تصمیم‌گیری انسانی همواره در شرایط \concept{عقلانیت
محدود}\footnote{\lat{Bounded Rationality}} صورت می‌گیرد%
\src{\lr{Simon, 1957, pp.\,196--206}}. این اصل در مورد نجات‌طلبی به‌شدت
صادق است: گروهی که در بحران ستم یا فروپاشی به سر می‌برد، ممکن است
اطلاعات ناقصی درباره‌ی نیات واقعی نیروی خارجی، هزینه‌های بلندمدت
مداخله، یا گزینه‌های جایگزین داشته باشد.


% ────────────────────────────────────────────────────────────────
\section{ادبیات نظری: سه مکتب و یک مسئله}
\label{sec:ch01-theories}

سه مکتب اصلی روابط بین‌الملل هر یک تفسیر متفاوتی از مداخله‌ی خارجی
ارائه می‌دهند:

\begin{tcolorbox}[mybox, title={\rl{سه خوانش نظری از مداخله‌ی خارجی}}]

\needspace{4\baselineskip}
\subsubsection{واقع‌گرایی}
\label{subsubsec:ch01-realism}
از منظر واقع‌گرایان --- \thinker{هانس مورگنتا}%
\footnote{\lat{Hans J.\,Morgenthau (1904--1980)}}
و \thinker{کنث والتز}%
\footnote{\lat{Kenneth N.\,Waltz (1924--2013)}} ---
مداخله‌ی خارجی صرفاً ابزاری در خدمت \concept{منافع ملی} و
\concept{موازنه‌ی قدرت} است. نجات‌بخشی لفاظی ایدئولوژیک است؛
آنچه اهمیت دارد محاسبه‌ی قدرت و امنیت
است\src{\lr{Morgenthau, 1948, ch.\,3; Waltz, 1979, ch.\,6}}.

\needspace{4\baselineskip}
\subsubsection{لیبرالیسم}
\label{subsubsec:ch01-liberalism}
لیبرال‌ها --- از \thinker{ایمانوئل کانت}%
\footnote{\lat{Immanuel Kant (1724--1804)}}
تا \thinker{مایکل دویل}%
\footnote{\lat{Michael W.\,Doyle (1948–)}} ---
معتقدند مداخله می‌تواند در خدمت \concept{صلح دموکراتیک} و
\concept{حقوق بشر} مشروع باشد، به‌ویژه هنگامی که یک رژیم
استبدادی حقوق بنیادین شهروندانش را نقض
می‌کند\src{\lr{Doyle, 1983, pp.\,205--235}}.

\needspace{4\baselineskip}
\subsubsection{پسااستعمارگرایی و نظریه‌ی انتقادی}
\label{subsubsec:ch01-postcolonial}
متفکران پسااستعماری --- از \thinker{ادوارد سعید}%
\footnote{\lat{Edward W.\,Said (1935--2003)}}
تا \thinker{فرانتس فانون}%
\footnote{\lat{Frantz Fanon (1925--1961)}} ---
هرگونه مداخله‌ی قدرت‌های بزرگ را ذاتاً مشکوک می‌دانند و آن را
بازتولید \concept{سلطه‌ی استعماری} در لباس جدید تفسیر
می‌کنند\src{\lr{Said, 1993, pp.\,xi--xxvi; Fanon, 1961, ch.\,1}}.

\end{tcolorbox}


% ────────────────────────────────────────────────────────────────
\section{هم‌چسبندگی و مداخله: کیش، آیین، ایدئولوژی، طبقه}
\label{sec:ch01-adhesion}

یکی از مهم‌ترین متغیرهای تعیین‌کننده در پویایی نجات‌طلبی آن است که
دعوت‌کننده و مداخله‌گر چه نوع \concept{هم‌چسبندگی}%
\footnote{منظور از هم‌چسبندگی، هرگونه پیوند هویتی غیرملی است که
  دعوت‌کننده را به نیروی خارجی متصل می‌کند.}
دارند. این هم‌چسبندگی‌ها را می‌توان در چند دسته طبقه‌بندی کرد:

\begin{enumerate}[label=\textbf{\alph*.}, nosep, rightmargin=1em]
  \item \concept{هم‌کیشی}: پیوند مذهبی-دینی (مسیحیان ارمنی و روسیه‌ی
    تزاری، شیعیان عراق و ایران)؛
  \item \concept{هم‌تباری}: پیوند قومی-نژادی (پان‌ترکیسم، پان‌اسلاویسم،
    پان‌ژرمنیسم)؛
  \item \concept{هم‌ایدئولوژی}: پیوند فکری-سیاسی (کمونیسم بین‌المللی،
    لیبرال-دموکراسی، اسلام‌گرایی)؛
  \item \concept{هم‌طبقه‌ای}: پیوند طبقاتی-اقتصادی (حمایت اتحاد شوروی
    از جنبش‌های کارگری، حمایت آمریکا از بورژوازی نوپا)؛
  \item \concept{هم‌منفعتی}: پیوند صرفاً ابزاری و فرصت‌طلبانه بدون
    ریشه‌ی هویتی عمیق.
\end{enumerate}

\begin{tcolorbox}[wavebox, title={\rl{نکته‌ی تحلیلی}}]
فرضیه‌ی محوری این پژوهش آن است که \textbf{شدت هم‌چسبندگی} با
\textbf{پایداری و نتایج مداخله} رابطه‌ی معناداری دارد: هم‌چسبندگی‌های
ارگانیک و تاریخی (مانند هم‌تباری یا هم‌کیشی دیرینه) الگوهای متفاوتی
از مداخله ایجاد می‌کنند تا هم‌چسبندگی‌های ابزاری و فرصت‌طلبانه.
این فرضیه در فصل ۱۲ آزمون تطبیقی خواهد شد.
\end{tcolorbox}


% ────────────────────────────────────────────────────────────────
\section{فعالیت‌های پیش‌زمینه‌ساز: میسیونری، ایدئولوژیک و استعماری}
\label{sec:ch01-precursors}

مداخله‌ی نظامی اغلب نقطه‌ی اوج یک فرایند بلندمدت‌تر است. پیش از رسیدن
به مرحله‌ی نظامی، معمولاً فعالیت‌هایی «نرم» زمینه را آماده کرده‌اند:

\subsection{فعالیت‌های میسیونری و تبلیغی}
\label{subsec:ch01-missionary}

مفهوم «میسیونری» در این پژوهش فراتر از معنای صرفاً مذهبی آن به
کار می‌رود. هر فعالیتی که با هدف \concept{تغییر باورها و ارزش‌های}
یک جامعه از بیرون صورت گیرد --- خواه مسیحی‌سازی، خواه صدور انقلاب
اسلامی، خواه ترویج لیبرال-دموکراسی، خواه پان‌ترکیسم ---
در این چارچوب «میسیونری» محسوب می‌شود.

\thinker{آنتونیو گرامشی}\footnote{\lat{Antonio Gramsci (1891--1937)},
فیلسوف و نظریه‌پرداز مارکسیست ایتالیایی.}
با مفهوم \concept{هژمونی فرهنگی}\footnote{\lat{Cultural Hegemony}}
نشان داده که سلطه پیش از آنکه نظامی باشد، فرهنگی
است\src{\lr{Gramsci, 1971, pp.\,12--13}}. فعالیت‌های میسیونری ---
در گسترده‌ترین معنای آن --- ابزار اصلی ساختن این هژمونی‌اند.

\subsection{استعمار اولیه و ثانویه}
\label{subsec:ch01-colonialism}

% ═══════════════════════════════════════════════════════════════
%  ch01.tex — ادامه (از تعریف استعمار اولیه و ثانویه)
% ═══════════════════════════════════════════════════════════════

\begin{tcolorbox}[defbox, title={تعریف: استعمار اولیه و ثانویه}]
\concept{استعمار اولیه}\footnote{\lat{Primary Colonialism}}
به فرایندی اطلاق می‌شود که در آن یک قدرت خارجی مستقیماً سرزمینی را
تصرف و اداره می‌کند --- مانند استعمار بریتانیا در هند یا فرانسه در
الجزایر. در این شکل، سلطه آشکار، نظامی و اداری است.

\concept{استعمار ثانویه}\footnote{\lat{Secondary / Neo-Colonialism}}
--- یا نو-استعمارگری --- به وضعیتی اطلاق می‌شود که در آن، قدرت
خارجی بدون حضور نظامی مستقیم و از طریق ابزارهای اقتصادی، فرهنگی
و دیپلماتیک سلطه‌ی خود را اعمال می‌کند. \thinker{کوامه نکرومه}%
\footnote{\lat{Kwame Nkrumah (1909--1972)},
نخست‌وزیر و رئیس‌جمهور غنا و نظریه‌پرداز نو-استعمارگری.}
در اثر خود با عنوان \lat{\textit{Neo-Colonialism: The Last Stage of
Imperialism}} (۱۹۶۵) استدلال می‌کند که نو-استعمار خطرناک‌تر از
استعمار کلاسیک است، زیرا «قدرت بدون مسئولیت» اعمال
می‌شود\src{\lr{Nkrumah, 1965, p.\,xi}}.
\end{tcolorbox}

پرسش کلیدی در اینجا آن است: آیا \concept{فعالیت‌های میسیونری و
ایدئولوژیک} را باید ذیل مقوله‌ی استعمار طبقه‌بندی کرد، یا آن‌ها
مقوله‌ی مستقلی‌اند؟ پاسخ این پرسش به‌شدت به بافت تاریخی وابسته است.
در برخی موارد --- مانند فعالیت‌های میسیونرهای مسیحی در آفریقای
قرن نوزدهم --- فعالیت تبلیغی و استعمار چنان درهم‌تنیده بودند که
تفکیکشان ناممکن
است\src{\lr{Comaroff \& Comaroff, 1991, pp.\,1--48}}.
اما در موارد دیگر --- مانند فعالیت‌های روشنفکران مشروطه‌خواه ایرانی
در قفقاز --- پیوند ایدئولوژیک بدون بُعد استعماری و حتی در تعارض با
منافع قدرت‌های استعماری شکل گرفته است.

\subsection{سه‌گانه‌ی تمایز: ارگانیک، ابزاری، امپریال}
\label{subsec:ch01-triad}

بر اساس بررسی ادبیات موجود، طبقه‌بندی سه‌گانه‌ی زیر را پیشنهاد
می‌کنیم:

\begin{figure}[htbp]
\centering
\begin{tikzpicture}[
  >=Stealth,
  every node/.style={font=\small},
  circlestyle/.style={
    circle, draw, minimum size=3.5cm,
    text width=2.8cm, align=center, font=\small\bfseries
  }
]
% سه دایره
\node[circlestyle, fill=BgTeal!60, text=AccentTeal] (org) at (0,0)
  {\rl{ارگانیک}\\[3pt]%
   {\footnotesize\rl{پیوند تاریخی}}\\%
   {\footnotesize\rl{هم‌تبار/هم‌کیش}}};

\node[circlestyle, fill=BgWarm!60, text=AccentAmber] (inst) at (5,0)
  {\rl{ابزاری}\\[3pt]%
   {\footnotesize\rl{محاسبه‌ی منفعت}}\\%
   {\footnotesize\rl{فرصت‌طلبانه}}};

\node[circlestyle, fill=red!10, text=AccentRed] (imp) at (2.5,4)
  {\rl{امپریال}\\[3pt]%
   {\footnotesize\rl{سلطه‌طلبی}}\\%
   {\footnotesize\rl{کشورگشایی}}};

% فلش‌ها
\draw[->, thick, NeutralDark] (org) -- (inst)
  node[midway, below, font=\scriptsize] {\rl{ابزاری شدن پیوند}};
\draw[->, thick, NeutralDark] (inst) -- (imp)
  node[midway, right, font=\scriptsize, text width=2cm, align=center]
  {\rl{تبدیل به سلطه}};
\draw[->, thick, NeutralDark, dashed] (org) -- (imp)
  node[midway, left, font=\scriptsize, text width=2cm, align=center]
  {\rl{پوشش هویتی}\\{\scriptsize\rl{برای سلطه}}};

% عنوان
\node[above, font=\normalsize\bfseries, text=PrimaryDeep]
  at (2.5,6.2) {\rl{سه‌گانه‌ی روابط مداخله‌ای}};
\end{tikzpicture}
\caption{سه‌گانه‌ی تمایز در روابط مداخله‌ای: ارگانیک، ابزاری و امپریال}
\label{fig:ch01-triad}
\end{figure}

\begin{enumerate}[label=\textbf{\arabic*.}, nosep, rightmargin=1em]
  \item \concept{مداخله‌ی ارگانیک}: میان جوامعی رخ می‌دهد که پیوند
    تاریخی، فرهنگی یا هویتی عمیق و دیرینه دارند. دعوت از «خویشاوند»
    صورت می‌گیرد نه از «بیگانه‌ی محض». نمونه: دعوت گرجی‌های مسیحی
    از روسیه‌ی ارتدوکس.
  \item \concept{مداخله‌ی ابزاری}: پیوند هویتی یا وجود ندارد یا
    ثانوی است. طرفین بر اساس محاسبه‌ی سود-زیان عمل می‌کنند. نمونه:
    اتحاد شیخ خزعل با بریتانیا.
  \item \concept{مداخله‌ی امپریال}: بازیگر خارجی اساساً به دنبال گسترش
    سلطه‌ی خویش است و «دعوت» یا «نجات‌بخشی» صرفاً بهانه‌ای برای
    مشروعیت‌بخشی است. نمونه: حمله‌ی شوروی به افغانستان (۱۹۷۹).
\end{enumerate}


% ────────────────────────────────────────────────────────────────
\section{معیارهای موفقیت و شکست: بازه‌ی زمانی و چندبعدی بودن}
\label{sec:ch01-success}

ارزیابی موفقیت یا شکست یک مداخله‌ی خارجی نمی‌تواند تک‌بعدی باشد.
\thinker{جیمز مریدیث}\footnote{\lat{James Meernik},
پژوهشگر علوم سیاسی دانشگاه شمال تگزاس.}
در بررسی ۳۵ مورد مداخله‌ی نظامی آمریکا (۱۹۴۵--۲۰۰۱) نشان داد که
«موفقیت» بسته به معیار مورد استفاده --- نظامی، سیاسی، اقتصادی یا
حقوق بشری --- نتایج کاملاً متفاوتی
دارد\src{\lr{Meernik, 2004, pp.\,67--89}}.

\subsection{ماتریس ارزیابی پیشنهادی}
\label{subsec:ch01-matrix}

\begin{table}[htbp]
\centering
\caption{ماتریس چندبعدی ارزیابی مداخلات خارجی}
\label{tab:ch01-matrix}
\begin{adjustbox}{max width=\textwidth}
\begin{tabular}{>{\bfseries\color{PrimaryMid}}p{2.5cm}
  p{3.5cm} p{3.5cm} p{3.5cm}}
\toprule
\textbf{بُعد} &
\textbf{کوتاه‌مدت (۰--۵ سال)} &
\textbf{میان‌مدت (۵--۲۵ سال)} &
\textbf{بلندمدت (بیش از ۲۵ سال)} \\
\midrule
نظامی &
پایان درگیری؟ &
ثبات امنیتی؟ &
بازتولید خشونت؟ \\
\midrule
سیاسی &
تغییر رژیم؟ &
نهادسازی دموکراتیک؟ &
تحکیم حاکمیت ملی؟ \\
\midrule
اقتصادی &
بازسازی اولیه؟ &
رشد اقتصادی پایدار؟ &
استقلال اقتصادی؟ \\
\midrule
حقوق بشری &
توقف نسل‌کشی/ستم؟ &
نهادینه‌شدن حقوق؟ &
فرهنگ حقوق بشر؟ \\
\midrule
فرهنگی-هویتی &
حفظ هویت محلی؟ &
تعامل فرهنگی سازنده؟ &
استقلال فرهنگی؟ \\
\midrule
روان‌شناختی &
احساس رهایی؟ &
اعتماد اجتماعی؟ &
غلبه بر ضربه‌ی روانی؟ \\
\bottomrule
\end{tabular}
\end{adjustbox}
\end{table}

\begin{tcolorbox}[wavebox, title={\rl{نکته‌ی روش‌شناختی}}]
در سرتاسر این پژوهش، هر مطالعه‌ی موردی بر اساس این ماتریس شش‌بعدی
و سه‌زمانی ارزیابی خواهد شد. این رویکرد امکان مقایسه‌ی تطبیقی
منسجم میان موارد بسیار متنوع --- از حمله‌ی اعراب به ایران ساسانی
تا مداخله‌ی ناتو در کوزوو --- را فراهم می‌آورد.
\end{tcolorbox}


% ────────────────────────────────────────────────────────────────
\section{عواقب و پیامدهای مداخلات: نگاهی کلی}
\label{sec:ch01-consequences}

پیامدهای مداخلات خارجی را می‌توان در چند مقوله‌ی کلی خلاصه کرد:

\begin{tcolorbox}[enemybox, title={\rl{پیامدهای بالقوه‌ی منفی}}]
\begin{enumerate}[label=\textbf{\arabic*.}, nosep, rightmargin=1em]
  \item \enemy{از دست رفتن استقلال}: مداخله‌گر پس از «نجات» حاضر به
    ترک نیست و سلطه‌ی خود را تثبیت
    می‌کند\src{\lr{Downes \& Monten, 2013, pp.\,94--132}}.
  \item \enemy{تحمیل دیدگاه‌ها}: ارزش‌ها و نهادهای بیگانه بدون
    توجه به بافت محلی تحمیل
    می‌شوند\src{\lr{Paris, 2004, pp.\,40--55}}.
  \item \enemy{استثمار منابع}: «نجات» بهانه‌ای برای دسترسی به منابع
    طبیعی و ثروت ملی می‌شود\src{\lr{Le Billon, 2012, pp.\,23--47}}.
  \item \enemy{هزینه‌های انسانی و مالی}: جنگ و مداخله خود منبع رنج
    و تخریب می‌شوند\src{\lr{Crawford, 2019, pp.\,1--15}}.
  \item \enemy{بی‌اعتمادی مزمن}: تجربه‌ی مداخله، اعتماد اجتماعی و
    بین‌المللی را برای نسل‌ها تخریب
    می‌کند\src{\lr{Fukuyama, 2004, pp.\,92--110}}.
\end{enumerate}
\end{tcolorbox}

\begin{tcolorbox}[wavebox, title={\rl{پیامدهای بالقوه‌ی مثبت}}]
\begin{enumerate}[label=\textbf{\arabic*.}, nosep, rightmargin=1em]
  \item \concept{توقف نسل‌کشی یا ستم سیستماتیک}: مانند مداخله‌ی
    ویتنام در کامبوج (۱۹۷۸) که به حکومت خمرهای سرخ پایان داد.
  \item \concept{گذار به دموکراسی}: مانند بازسازی ژاپن و آلمان
    پس از جنگ جهانی دوم\src{\lr{Dower, 1999; Judt, 2005}}.
  \item \concept{نهادسازی}: ایجاد نهادهای حقوقی و اداری مدرن
    در جوامعی که فاقد آن بودند.
  \item \concept{بهبود ساختاری}: اصلاحات بهداشتی، آموزشی و
    زیرساختی --- حتی اگر با انگیزه‌ی استعماری صورت گرفته باشند.
\end{enumerate}
\end{tcolorbox}


% ────────────────────────────────────────────────────────────────
\section{جمع‌بندی فصل و نقشه‌ی راه}
\label{sec:ch01-summary}

در این فصل، مفاهیم بنیادین پژوهش حاضر --- مداخله‌ی خارجی، نجات‌بخشی،
نجات‌طلبی، هم‌چسبندگی، و طیف نرم-سخت مداخله --- تبیین شد. سه
طبقه‌بندی نظری اصلی (واقع‌گرایی، لیبرالیسم، پسااستعمارگرایی) و
یک سه‌گانه‌ی تحلیلی (ارگانیک، ابزاری، امپریال) معرفی گردید. همچنین
ماتریس ارزیابی شش‌بعدی و سه‌زمانی به‌عنوان ابزار تحلیلی اصلی
پژوهش ارائه شد.

در فصل بعد، به مهم‌ترین و پیچیده‌ترین پرسش این پژوهش خواهیم پرداخت:
\textbf{رابطه‌ی دوسویه و چندسویه‌ی نجات‌بخشی و نجات‌طلبی} و
تنش‌های بنیادینی که این رابطه تولید می‌کند --- از پارادوکس اجبار
به آزادی تا مسئله‌ی بلوغ درون‌زاد ارزش‌ها.